\documentclass{beamer}
\usepackage[utf8]{inputenc}

\usetheme{Madrid}
\usecolortheme{default}
\usepackage{amsmath,amssymb,amsfonts,amsthm}
\usepackage{txfonts}
\usepackage{tkz-euclide}
\usepackage{listings}
\usepackage{adjustbox}
\usepackage{array}
\usepackage{tabularx}
\usepackage{gvv}
\usepackage{lmodern}
\usepackage{circuitikz}
\usepackage{tikz}
\usepackage{graphicx}

\setbeamertemplate{page number in head/foot}[totalframenumber]

\usepackage{tcolorbox}
\tcbuselibrary{minted,breakable,xparse,skins}



\definecolor{bg}{gray}{0.95}
\DeclareTCBListing{mintedbox}{O{}m!O{}}{%
	breakable=true,
	listing engine=minted,
	listing only,
	minted language=#2,
	minted style=default,
	minted options={%
		linenos,
		gobble=0,
		breaklines=true,
		breakafter=,,
		fontsize=\small,
		numbersep=8pt,
		#1},
	boxsep=0pt,
	left skip=0pt,
	right skip=0pt,
	left=25pt,
	right=0pt,
	top=3pt,
	bottom=3pt,
	arc=5pt,
	leftrule=0pt,
	rightrule=0pt,
	bottomrule=2pt,
	toprule=2pt,
	colback=bg,
	colframe=orange!70,
	enhanced,
	overlay={%
		\begin{tcbclipinterior}
			\fill[orange!20!white] (frame.south west) rectangle ([xshift=20pt]frame.north west);
	\end{tcbclipinterior}},
	#3,
}
\lstset{
	language=C,
	basicstyle=\ttfamily\small,
	keywordstyle=\color{blue},
	stringstyle=\color{orange},
	commentstyle=\color{green!60!black},
	numbers=left,
	numberstyle=\tiny\color{gray},
	breaklines=true,
	showstringspaces=false,
}
%------------------------------------------------------------
%This block of code defines the information to appear in the
%Title page
\title %optional
{12.684}
%\subtitle{A short story}

\author % (optional)
{RAVULA SHASHANK REDDY - EE25BTECH11047}

 \begin{document}
	
	
	\frame{\titlepage}
	\begin{frame}{Question}
    If $\vec{A}$ is a square matrix then the orthogonality property mandates:

\begin{enumerate}[label=\alph*)]
\item $\vec{A}\vec{A}^\top = \vec{I}$
\item $\vec{A}\vec{A}^\top = 0$
\item $\vec{A}\vec{A}^\top = \vec{A}^{-1}$
\item $\vec{A}\vec{A}^\top = \vec{A}$
\end{enumerate}
\end{frame}
\begin{frame}{Solution}
    A square matrix $\vec{A}$ is orthogonal if its transpose equals its inverse: 
\begin{align}
& \vec{A}^\top = \vec{A}^{-1} 
\end{align}

Multiplying both sides on the left by $\vec{A}$:
\begin{align}
& \vec{A} \vec{A}^\top = \vec{A} \vec{A}^{-1} \notag \\
& \implies \vec{A}\vec{A}^\top = \vec{I}
\end{align}
\end{frame}
\begin{frame}[fragile]
\frametitle{C Code}
\begin{lstlisting}
    #include <stdio.h>
#include <math.h>

#define N 3  // Dimension of the square matrix
#define EPS 1e-6  // Tolerance for floating-point comparison

int main() {
    int i, j, k;
    double A[N][N] = {
        {1/sqrt(2), -1/sqrt(2), 0},
        {1/sqrt(2),  1/sqrt(2), 0},
        {0,          0,         1}
    };
    double AA_T[N][N] = {0};

    // Compute AA^T
    for(i = 0; i < N; i++) {
        for(j = 0; j < N; j++) {
        \end{lstlisting}
        \end{frame}
\begin{frame}[fragile]
\frametitle{C Code}
\begin{lstlisting}
            for(k = 0; k < N; k++) {
                AA_T[i][j] += A[i][k] * A[j][k]; 
            }
        }
    }
    int isOrthogonal = 1;
    for(i = 0; i < N; i++) {
        for(j = 0; j < N; j++) {
            if(i == j && fabs(AA_T[i][j] - 1) > EPS) isOrthogonal = 0;
            if(i != j && fabs(AA_T[i][j]) > EPS) isOrthogonal = 0;
        }
    }
if(isOrthogonal)printf("Matrix A is orthogonal");
else
printf("Matrix A is not orthogonal");
    return 0;
}
\end{lstlisting}
\end{frame}
\begin{frame}[fragile]
\frametitle{Python Direct}
\begin{lstlisting}
    import numpy as np

# Define a square matrix A
A = np.array([
    [1/np.sqrt(2), -1/np.sqrt(2), 0],
    [1/np.sqrt(2),  1/np.sqrt(2), 0],
    [0,             0,            1]
])

# Compute A * A^T
AA_T = A @ A.T

# Check if AA^T is identity
if np.allclose(AA_T, np.eye(A.shape[0])):
    print("Matrix A is orthogonal.")
else:
    print("Matrix A is not orthogonal.")
\end{lstlisting}
\end{frame}

\end{document}
